% ================================================================
% Section: Deployment on Resource-Constrained Environments
% From main (2).tex — Section VI
% ================================================================

\section{Deployment on Resource-Constrained Environments}
\label{sec:deployment}

% ----------------------------------------------------------------
\subsection{On-Device and Mobile Deployment}
\label{subsec:deployment_mobile}

On-device deployment eliminates the latency, cost, and privacy exposure of cloud API inference. It enables AI functionality in offline and low-bandwidth environments and is a prerequisite for truly private personal assistants. The primary challenges are memory (RAM and storage), sustained inference throughput under thermal constraints, and battery impact.

\subsubsection{Apple Intelligence (AFM)}

Apple's Foundation Models \citep{apple2024intelligence,gunter2025afm} are the most fully documented on-device \gls{slm} deployment. The 2025 report \citep{gunter2025afm} describes a $\sim$3B on-device model, optimized for Apple silicon through KV-cache sharing (dividing the model so the second block reuses the first block's cache, a reported 37.5\% memory saving) and 2-bit quantization-aware training, paired with a Parallel-Track Mixture-of-Experts server model on Private Cloud Compute. The on-device model is reported competitive with Qwen3-4B and Gemma-3-4B on its target tasks; the server model with Llama~4 Scout. The third-generation framework \citep{apple2026afm3} adds a natively multimodal, sparsely activated on-device model and opens the same API surface to third-party cloud models, so that an application calls the local \gls{slm} by default and escalates to a frontier \gls{llm} only on demand. Google's Gemini Nano and Samsung's Gauss \citep{samsung2023gauss} occupy the same design point on Android and Galaxy devices.

\subsubsection{Samsung Gauss}

Samsung Gauss \citep{samsung2023gauss} is a family of language models for Galaxy devices, including a ${\sim}1.7$B on-device model targeting the Gauss1 chip and Exynos 2400 \gls{npu}. Samsung reports that Gauss enables on-device email summarization, translation, and writing assistance without requiring internet connectivity.

\subsubsection{General Mobile Deployment Considerations}

\glspl{slm} targeting mobile deployment must navigate three constraints. (1) \emph{Memory}: Consumer smartphones typically provide 6--12~GB RAM, of which the OS uses 2--4~GB, leaving 4--8~GB for applications. A 3B model at Q4 requires ${\sim}2$~GB, comfortable for most flagship devices. (2)~\emph{Thermal}: Sustained inference on mobile \glspl{npu} is throttled by thermals; realistic sustained throughput is lower than peak throughput benchmarks suggest. (3)~\emph{Battery}: vendor guidance for mobile-targeted \glspl{slm} places per-prompt battery cost low enough to permit hundreds of interactions per charge, though published per-device figures remain scarce and are not reported under a common protocol. The fundamental constraint in edge deployment is memory bandwidth rather than compute: flagship mobile and laptop devices offer roughly 50--90~GB/s of memory bandwidth (single-board computers well under 30~GB/s) versus 2--3~TB/s for datacenter GPUs \citep{xu2024device}, making INT4 quantization not merely a storage optimization but a critical performance enabler.

% ----------------------------------------------------------------
\subsection{Edge and IoT Deployment}
\label{subsec:deployment_edge}

Below smartphone-class hardware lies a diverse ecosystem of edge devices: Raspberry Pi single-board computers, NVIDIA Jetson modules, and microcontrollers. \glspl{slm} targeting this tier typically require:

\begin{description}
    \item[Sub-1B or extreme quantization] A Raspberry Pi~5 (8~GB RAM, Cortex-A76 CPU) runs Q4-quantized 1B models at roughly ten tokens per second, with sub-500M models comfortably faster. Benchmarking shows that llamafile achieves up to $4\times$ higher throughput and 30--40\% lower power consumption than Ollama for equivalent models on single-board computers \citep{costa2025sbc}.

    \item[Jetson Orin] NVIDIA's Jetson Orin modules with TensorRT-LLM \citep{nvidia_tensorrt2024} serve quantized 3 to 7B models at interactive rates, with the larger-memory NX and AGX parts faster than the Orin Nano, enabling real-time local inference for industrial and robotics applications.

    \item[Microcontrollers (MCU)] Sub-100M models using aggressive quantization (INT4/INT2) can run on ARM Cortex-M series MCUs with 512~KB--2~MB SRAM for specialized tasks like keyword detection and sensor data analysis.
\end{description}

Table~\ref{tab:appendix_hardware} details the specifications of the four
representative devices referenced throughout this section, from the Raspberry~Pi
5 up to a laptop-class CPU.

\begin{table}[ht]
\centering
\caption{Detailed specifications of the devices referenced in this section.}
\label{tab:appendix_hardware}
\footnotesize\adjustbox{max width=\linewidth}{%
\begin{tabular}{lcccc}
\toprule
\textbf{Specification} & \textbf{Raspberry Pi 5} & \textbf{Jetson Orin Nano} & \textbf{Galaxy S24+} & \textbf{Laptop} \\
\midrule
CPU & BCM2712 4$\times$A76 @ 2.4GHz & 6$\times$A78AE @ 1.5GHz & Snapdragon 8 Gen 3 & Ryzen 7 7840HS \\
GPU & VideoCore VII & 1024 CUDA (Ampere) & Adreno 750 & Radeon 780M \\
NPU & None & None & Hexagon & Ryzen AI \\
RAM & 8 GB LPDDR4X & 8 GB LPDDR5 (unified) & 12 GB LPDDR5X & 32 GB DDR5-5600 \\
Memory bandwidth & $\sim$17 GB/s & 68 GB/s & 51.2 GB/s & 89.6 GB/s \\
Storage & microSD / NVMe & NVMe & UFS 4.0 & NVMe SSD \\
TDP & 12W & 25W & $\sim$8W sustained & 35--54W \\
Inference engine & llama.cpp & MLC LLM & ExecuTorch & llama.cpp \\
\bottomrule
\end{tabular}%
}
\end{table}

% ----------------------------------------------------------------
\subsection{Deployment Frameworks}
\label{subsec:deployment_frameworks}

A mature ecosystem of open-source frameworks enables \gls{slm} deployment across hardware targets. Table~\ref{tab:deployment_frameworks} provides a structured comparison.

\begin{table}[ht]
\caption{Key SLM deployment frameworks, their primary targets, and characteristics. Throughput figures in the Notes column are vendor- or project-reported on the hardware each names and are not comparable across rows.}
\label{tab:deployment_frameworks}
\centering
\footnotesize\adjustbox{max width=\linewidth}{%
\begin{tabular}{@{}lllp{5.5cm}@{}}
\toprule
\textbf{Framework} & \textbf{Primary Target} & \textbf{Quantization} & \textbf{Notes} \\
\midrule
llama.cpp \citep{gerganov2023llamacpp}   & CPU, Metal, CUDA, ROCm & GGUF (Q2--Q8)   & De facto standard; minimal dependencies; ecosystem includes Ollama, llamafile \\
ONNX Runtime \citep{microsoft2019onnx} & GPU, CPU, NPU, WASM & INT4/INT8       & Cross-platform; DirectML/CoreML/CUDA backends; Phi-3-mini at 35~tok/s on Snapdragon NPU \\
MLX \citep{mlx2023}             & Apple Silicon          & 4-bit           & Unified memory; eliminates GPU transfers; 55+ tok/s for LLaMA-3.2-3B on M3~Pro \\
TensorRT-LLM \citep{nvidia_tensorrt2024} & NVIDIA GPUs   & INT4/INT8/FP16  & Continuous batching; kernel fusion; 2--3$\times$ over PyTorch for 7B models \\
ExecuTorch \citep{meta2024executorch} & Mobile (iOS, Android) & INT4          & 50~KB base runtime; 12+ hardware backends; deployed in Instagram, WhatsApp \\
MLC-LLM \citep{mlc2024}        & iOS, Android, web, MCU & 4-bit          & TVM compilation; WebGPU support; universal platform coverage \\
WebLLM \citep{webllm2024}       & Web browsers           & 4-bit WASM     & WebGPU acceleration; 80\% of native performance; OpenAI-compatible API \\
\bottomrule
\end{tabular}%
}
\end{table}

% ----------------------------------------------------------------
\subsection{Why the Deployment Gap Matters}
\label{subsec:deployment_latency}

The practical consequence of the preceding sections is a gap that is structural
rather than marginal. A quantized 3B model occupies roughly 2~GB of weight
memory and runs on hardware the user already owns, so its marginal cost per
query is zero and no data leaves the device. A 70B model at FP16 needs roughly
140~GB, which places it on multi-accelerator server hardware, and a frontier API
model adds a network round trip and a per-token charge to every call. Between
those two conditions there is no continuum: the small model is either adequate
for the task or it is not, and where it is adequate the large model buys nothing
that the task can use.

Published cost estimates put the gap at one to two orders of magnitude.
\citet{belcak2025slm_agentic} estimate that serving a 7B \gls{slm} is 10 to
30$\times$ cheaper than serving a 70 to 175B \gls{llm} once latency, energy, and
FLOPs are accounted for. On the latency side, the binding constraint on edge
hardware is prefill rather than generation, because time-to-first-token
dominates perceived responsiveness in interactive use, and prefill cost scales
with context length. This is why \gls{gqa} and local-global attention
interleaving, which cut \gls{kv}-cache traffic, matter more for on-device
deployment than for server deployment.

We do not tabulate per-scenario latency and cost figures here. Serving prices
change faster than a survey's publication cycle, and published throughput
figures are not measured under comparable batch sizes, context lengths, or
backends, so a table of them would imply a precision the underlying reporting
does not support. The relevant claim for the argument of
Section~\ref{sec:agentic} is ordinal and robust to that imprecision: for a
bounded, repetitive call that a small model can complete correctly, the small
model is cheaper by a wide margin, lower in latency once the network hop is
counted, and available offline.
